\section{Projektbeschreibung}
\label{sec:project}

\subsection{Ziel und Systemidee}

Ziel des Projekts ist ein Demonstrator für ein geschwindigkeitsabhängiges
Schutzsystem eines rotierenden Aktors. Das Prinzip orientiert sich vereinfacht
an \emph{Speed and Separation Monitoring} (SSM): Mit steigender
Geschwindigkeitsstufe wird ein größerer Sicherheitsabstand verwendet.
Unterschreitet ein Objekt diesen Abstand, wird die Motoransteuerung abgeschaltet.

Der Demonstrator dient ausschließlich zur Untersuchung eines
sicherheitsgerichteten Prinzips. Er ist kein zertifiziertes Safety-System.
Insbesondere werden Raspberry Pi 4B, HC-SR04, Bluetooth und der verwendete
Motortreiber nicht als Safety-zertifizierte Komponenten betrachtet.

\keybox{Sicherer Zustand: Motor AUS -- der Rotor wird nicht aktiv angetrieben.}

Ein Software-Stopp bedeutet nicht automatisch mechanischen Stillstand. Der Rotor
kann aufgrund seiner Massenträgheit nachlaufen. Diese Nachlaufzeit muss bei der
Festlegung der Sicherheitsabstände berücksichtigt werden.

\subsection{Systemaufbau}

Das System besteht aus zwei Raspberry Pi 4B, zwei Ultraschall-Abstandssensoren
HC-SR04, einem Elektromotor mit Rotor, einem Motortreiber, einem LCD-Display und
einer Bluetooth-Verbindung zwischen den Raspberry Pis.

\begin{table}[H]
    \centering
    \small
    \begin{tabularx}{\textwidth}{@{}p{3.7cm}YY@{}}
        \toprule
        \textbf{Komponente} & \textbf{Aufgabe} & \textbf{Bedeutung für Safety} \\
        \midrule
        Raspberry Pi 1 & Bedienung, Auswahl der Geschwindigkeitsstufe, LCD-Anzeige, Logging und Bluetooth-Kommunikation & Liefert Sollwerte und zeigt Zustände an. LCD und Logging sind keine Safety-Funktionen. \\
        Raspberry Pi 2 & Sensorerfassung, Plausibilitätsprüfung, Safety-Logik und Motorsteuerung & Führt die lokale Safety-Entscheidung aus und kann den Motor unabhängig von Pi 1 abschalten. \\
        2 $\times$ HC-SR04 & Abstandsmessung zum Gefahrenbereich & Zwei gleiche Sensoren ermöglichen einen Quervergleich, bilden aber keine nachgewiesene Safety-Redundanz. \\
        Motortreiber / Motor & Umsetzung der Motoransteuerung und rotierende Gefahrenquelle & Motor AUS entspricht PWM = 0 und deaktiviertem Enable-Signal. Mechanischer Stillstand ist dadurch nicht garantiert. \\
        Bluetooth & Bidirektionale Kommunikation zwischen Pi 1 und Pi 2 & Wird mit Heartbeat, Timeout und Plausibilitätsprüfungen überwacht. \\
        \bottomrule
    \end{tabularx}
    \caption{Komponenten und Aufgaben des Demonstrators}
\end{table}

\subsection{Funktionsablauf}

Nach dem Start initialisieren beide Raspberry Pis ihre benötigten Schnittstellen
und Threads. Raspberry Pi 2 hält die Motoransteuerung während der
Initialisierung im sicheren Zustand. Erst wenn gültige Sensordaten, eine gültige
Kommunikation und eine gültige Geschwindigkeitsstufe vorliegen, darf der Motor
freigegeben werden.

Raspberry Pi 1 sendet die gewünschte Geschwindigkeitsstufe
\texttt{LOW}, \texttt{MEDIUM} oder \texttt{HIGH} sowie einen Heartbeat an
Raspberry Pi 2. Raspberry Pi 2 liest beide Abstandssensoren aus, prüft die
Messwerte und wählt den zugehörigen Sicherheitsabstand. Befindet sich ein Objekt
innerhalb dieses Abstands oder schlägt eine Safety-Prüfung fehl, wird der Motor
abgeschaltet.

Raspberry Pi 2 sendet Status, Abstandswerte und Fehlermeldungen an Raspberry Pi
1 zurück. Diese Informationen werden am LCD angezeigt und protokolliert.

Nach einer Abschaltung läuft der Motor ohne Quittierung wieder an, sobald alle
Prüfungen bestanden sind (Wiederanlaufbedingung, siehe
Abschnitt~\ref{sec:states}).

\subsection{Zentrale Designparameter}

Die folgenden Werte sind Projektparameter für den Demonstrator und müssen am
realen Aufbau getestet und gegebenenfalls angepasst werden.

\begin{table}[H]
    \centering
    \begin{tabularx}{\textwidth}{@{}p{5.1cm}p{3.2cm}Y@{}}
        \toprule
        \textbf{Parameter} & \textbf{Startwert} & \textbf{Bedeutung} \\
        \midrule
        Heartbeat Pi 1 $\rightarrow$ Pi 2 & 100 ms (10 Hz) & Zyklische Lebensmeldung und Übertragung der Sollstufe \\
        Kommunikations-Timeout & 300 ms & Keine gültige Nachricht innerhalb dieses Zeitraums führt zu Motor AUS \\
        Safety-Logik-Zyklus & 50 ms & Zyklische Prüfung der Safety-Bedingungen auf Pi 2 \\
        Max. Alter der Sensordaten & 150 ms & Ältere Messwerte dürfen nicht für eine Freigabe verwendet werden \\
        Gültiger Messbereich & 2--400 cm & Werte außerhalb des Messbereichs des HC-SR04 sind ungültig \\
        Max. Abstandsänderung & 0{,}2 cm/ms (= 2 m/s) & Schnellere Änderung zwischen zwei Messungen eines Sensors ist unplausibel \\
        Sensorabweichung & $> 5$ cm in zwei aufeinanderfolgenden Messzyklen & Sensorwerte werden als unplausibel bewertet \\
        Sicherheitsabstände & $S_{\mathrm{LOW}} < S_{\mathrm{MEDIUM}} < S_{\mathrm{HIGH}}$ & Endgültige Werte werden am Demonstrator anhand von Reaktions- und Nachlaufzeit festgelegt \\
        \bottomrule
    \end{tabularx}
    \caption{Zentrale Startwerte des Designs}
\end{table}

\subsection{Grenzen des Demonstrators}
\label{sec:limits}

Die tatsächliche Motordrehzahl wird nicht mit einem Encoder gemessen. Daher darf
nicht angenommen werden, dass ein PWM-Wert proportional zur realen Drehzahl ist.
Die Geschwindigkeitsstufen stellen lediglich definierte Ansteuerungsklassen dar.
Außerdem können zwei gleiche HC-SR04 durch eine gemeinsame Fehlerursache
ähnliche falsche Werte liefern. Ein Defekt des Motortreibers, bei dem der Motor
trotz Abschaltsignal weiter bestromt wird, kann mit der vorhandenen Architektur
nicht sicher beherrscht werden.

Der Wiederanlauf erfolgt automatisch, ohne Quittierung durch den Bediener.
Nach einem Fehler oder einem Objekt-Stopp kann der Motor daher wieder anlaufen,
ohne dass der Bediener dies ausdrücklich freigegeben hat. Befindet sich ein
Objekt genau an der Grenze des Sicherheitsabstands, kann der Motor
wiederholt ein- und ausschalten. Diese Einschränkungen werden für den
Demonstrator bewusst akzeptiert. In einer industriellen Anwendung wäre ein
unerwarteter Wiederanlauf durch eine Quittierung und eine Hysterese zu
verhindern.
